\medskip
\noindent\textup{(A)}
设 $k$ 为特征 $0$ 的代数闭域, $S$ 为 $\operatorname{Spec}(k[T])$ 在 $0$ 处的 Hensel 化. 我们在 $b(x)$ 为一致化元时证明 3.3.5.

由 Lefschetz 原理, 归结到 $k=\mathbb C$ 的情形. 由于问题是局部的, 可假设 $f:X\to S$ 由映射
\[
\sum_i z_i^2:\mathbb C^{n+1}\longrightarrow\mathbb C.
\]
经基变换得到. 上面给出的 $\pm\delta$ 的构造在整系数上同调中有超越的类似物. 因而通过 XIV, 2.1, 类 $\pm\delta$ 对应于一个超越消灭圈, 即
\[
\pm\delta_{\mathrm{cl}}
\in H^n_{\{0\}}\bigl(X_0,R\Psi_{\mathrm{cl}\eta}(\mathbb Z)\bigr).
\]
的像. 由 XIV, 2.1 以及代数圈 $\pm\delta$ 的相应性质, $\pm\delta_{\mathrm{cl}}$ 在
\[
H^n_{\{0\}}\bigl(X_0,R\Psi_{\mathrm{cl}\eta}(\mathbb Z)\bigr)
\otimes\mathbb Z_\ell
\]
中的像为此群的生成元. 因此 $\pm\delta_{\mathrm{cl}}$ 为 $H^n_{\{0\}}(X_0,R\Psi_{\mathrm{cl}\eta}(\mathbb Z))$ 的两个生成元, 并与 XIV, 3.2 中的消灭圈一致. 公式 3.3.5 因而由 XIV, 3.2.11 得出. 严格来说, 应当检验 XIV, 2.1 的同构与杯积及迹的相容性.

\medskip
\noindent\textup{(B)}
我们对 (A) 中的 $X/S$ 证明 3.3.6, 再假设 $X/S$ 固有. 这由 XIII, 2.4.6.2 得出.

\medskip
\noindent\textup{(C)}
设 $p$ 为素数, $S_0$ 为 $\operatorname{Spec}(\mathbb Z[T])$ 在 $(p,T)$ 处的严格 Hensel 化. 我们对满足 $b(x)=T$ 的 $X/S_0$ 证明 3.3.6.

层 $R^if_*\Lambda$ 在 $\operatorname{Spec}(\mathbb Z[T])$ 在 $(T)$ 处的严格 Hensel 化上的限制由 (B) 控制. 此层在光滑除子 $T=0$ 的补集上局部常值, 而此除子具有混合特征. 因此断言 (C) 由 (B) 和 Abhyankar 引理得出.